% geometry/curvature_transitions.tex

\chapter{Curvature Transitions}

% ============================================================
\section{Motivation}
% ============================================================

\begin{remark}
Railway alignments rarely consist of isolated constant-curvature
segments.
\end{remark}

\begin{remark}
Instead, operational geometry requires smooth transitions between
different curvature states.
\end{remark}

\begin{remark}
Within AIM, transition curves are interpreted fundamentally as
curvature-transition processes.
\end{remark}

\begin{remark}
This shifts the interpretation:
\begin{itemize}
\item from position interpolation,
\item to controlled curvature evolution.
\end{itemize}
\end{remark}

% ============================================================
\section{Transition Principle}
% ============================================================

\begin{definition}[Curvature transition]
A curvature transition is a curvature law
\[
\kappa(s)
\]
connecting two curvature states:
\[
\kappa_A
\rightarrow
\kappa_B.
\]
\end{definition}

% ------------------------------------------------------------

\begin{remark}
The transition determines:
\begin{itemize}
\item geometric smoothness,
\item dynamic behaviour,
\item realization stability,
\item and operational comfort.
\end{itemize}
\end{remark}

% ------------------------------------------------------------

\begin{principle}[Transition principle]
Within AIM, transition geometry is governed primarily by curvature
evolution rather than by coordinate interpolation.
\end{principle}

% ============================================================
\section{Normalized Transition Space}
% ============================================================

\begin{definition}[Normalized transition parameter]
Transitions are commonly normalized to:
\[
w\in[0,1].
\]
\end{definition}

% ------------------------------------------------------------

\begin{definition}[Normalized curvature law]
A normalized transition law satisfies:
\[
\hat\kappa(w)\in[0,1].
\]
\end{definition}

% ------------------------------------------------------------

\begin{remark}
Normalization separates:
\begin{itemize}
\item intrinsic transition shape,
\item from physical scaling.
\end{itemize}
\end{remark}

% ------------------------------------------------------------

\begin{remark}
This enables:
\begin{itemize}
\item reusable transition families,
\item family interpolation,
\item lookup-based evaluation,
\item and normalized optimization.
\end{itemize}
\end{remark}

% ============================================================
\section{Transition Families}
% ============================================================

\begin{definition}[Transition family]
A transition family is a class of normalized curvature laws sharing
common structural properties.
\end{definition}

% ------------------------------------------------------------

\begin{remark}
Examples include:
\begin{itemize}
\item clothoids,
\item polynomial transitions,
\item trigonometric transitions,
\item Bloss transitions,
\item Vienna transitions,
\item and hybrid transition families.
\end{itemize}
\end{remark}

% ============================================================
\section{Clothoid Transition}
% ============================================================

\begin{definition}[Clothoid]
The clothoid is defined by:
\[
\kappa(s)
=
a s.
\]
\end{definition}

% ------------------------------------------------------------

\begin{remark}
The clothoid provides linear curvature evolution.
\end{remark}

% ------------------------------------------------------------

\begin{remark}
Historically, the clothoid became dominant because:
\begin{itemize}
\item it is simple,
\item analytically tractable,
\item and dynamically smoother than curvature jumps.
\end{itemize}
\end{remark}

% ============================================================
\section{Polynomial Transitions}
% ============================================================

\begin{definition}[Polynomial transition]
A polynomial transition satisfies:
\[
\kappa(s)
=
\sum_{i=0}^n a_i s^i.
\]
\end{definition}

% ------------------------------------------------------------

\begin{remark}
Polynomial transitions allow:
\begin{itemize}
\item higher continuity,
\item endpoint constraints,
\item and optimized dynamic behaviour.
\end{itemize}
\end{remark}

% ============================================================
\section{Trigonometric Transitions}
% ============================================================

\begin{definition}[Trigonometric transition]
A trigonometric transition contains sinusoidal curvature components.
\end{definition}

% ------------------------------------------------------------

\begin{remark}
These transitions often provide smoother higher-order behaviour than
low-degree polynomial transitions.
\end{remark}

% ============================================================
\section{Half-Wave Interpretation}
% ============================================================

\begin{definition}[Half-wave transition]
A half-wave transition is a normalized curvature segment connecting:
\[
0 \rightarrow 1
\]
or
\[
1 \rightarrow 0.
\]
\end{definition}

% ------------------------------------------------------------

\begin{remark}
Within AIM, many transition families are interpreted compositionally as:
\[
\text{halfWave}_1
\rightarrow
\text{core}
\rightarrow
\text{halfWave}_2.
\]
\end{remark}

% ------------------------------------------------------------

\begin{remark}
This enables:
\begin{itemize}
\item modular transition construction,
\item asymmetric transitions,
\item and family interpolation.
\end{itemize}
\end{remark}

% ============================================================
\section{Curvature Anchors}
% ============================================================

\begin{definition}[Curvature anchors]
Normalized transition families may define anchor values:
\[
0,
\quad
c_1,
\quad
c_2,
\quad
1.
\]
\end{definition}

% ------------------------------------------------------------

\begin{remark}
Anchors describe:
\begin{itemize}
\item transition partitioning,
\item continuity conditions,
\item and normalization structure.
\end{itemize}
\end{remark}

% ============================================================
\section{Transition Continuity}
% ============================================================

\begin{definition}[Transition continuity]
Transition quality depends on continuity of:
\[
\kappa(s),
\qquad
\kappa'(s),
\qquad
\kappa''(s).
\]
\end{definition}

% ------------------------------------------------------------

\begin{remark}
Higher continuity improves:
\begin{itemize}
\item jerk behaviour,
\item realization smoothness,
\item runtime stability,
\item and passenger comfort.
\end{itemize}
\end{remark}

% ============================================================
\section{Sparse Transition Representation}
% ============================================================

\begin{definition}[Sparse transition]
A sparse transition representation stores:
\begin{itemize}
\item transition family identity,
\item normalization parameters,
\item scaling parameters,
\item and reconstruction rules.
\end{itemize}
\end{definition}

% ------------------------------------------------------------

\begin{remark}
Within AIM, sparse transitions are interpreted as compact curvature
operators rather than dense geometric curves.
\end{remark}

% ============================================================
\section{Lookup-Based Transitions}
% ============================================================

\begin{definition}[Lookup-based transition]
A lookup-based transition uses:
\begin{itemize}
\item sampled normalized curvature,
\item interpolation rules,
\item and runtime reconstruction.
\end{itemize}
\end{definition}

% ------------------------------------------------------------

\begin{remark}
This enables:
\begin{itemize}
\item arbitrary transition families,
\item experimentally designed transitions,
\item realization-adapted transitions,
\item and runtime-efficient evaluation.
\end{itemize}
\end{remark}

% ============================================================
\section{Transition Derivatives}
% ============================================================

\begin{definition}[Transition derivatives]
Important quantities include:
\[
\kappa'(s),
\qquad
\kappa''(s),
\qquad
\kappa'''(s).
\]
\end{definition}

% ------------------------------------------------------------

\begin{remark}
These derivatives influence:
\begin{itemize}
\item jerk,
\item dynamic smoothness,
\item realization behaviour,
\item and optimization stability.
\end{itemize}
\end{remark}

% ============================================================
\section{Transition Symmetry}
% ============================================================

\begin{definition}[Symmetric transition]
A transition is symmetric if:
\[
\kappa(s)
=
\kappa(L-s).
\]
\end{definition}

% ------------------------------------------------------------

\begin{definition}[Asymmetric transition]
A transition is asymmetric if this symmetry is intentionally violated.
\end{definition}

% ------------------------------------------------------------

\begin{remark}
Asymmetric transitions may improve:
\begin{itemize}
\item realization quality,
\item vehicle dynamics,
\item or geometric fitting.
\end{itemize}
\end{remark}

% ============================================================
\section{Transition Frequency Behaviour}
% ============================================================

\begin{remark}
Transitions may be interpreted spectrally through their curvature
content.
\end{remark}

% ------------------------------------------------------------

\begin{remark}
High-frequency curvature behaviour influences:
\begin{itemize}
\item realization sensitivity,
\item machine smoothing,
\item and dynamic roughness.
\end{itemize}
\end{remark}

% ------------------------------------------------------------

\begin{proposition}[Spectral transition principle]
Transition quality depends not only on curvature continuity, but also on
spectral behaviour within curvature space.
\end{proposition}

% ============================================================
\section{Transition Optimization}
% ============================================================

\begin{remark}
Within AIM, optimization acts primarily on:
\begin{itemize}
\item transition families,
\item normalization structure,
\item transition lengths,
\item and curvature evolution.
\end{itemize}
\end{remark}

% ------------------------------------------------------------

\begin{remark}
Thus, optimization is interpreted primarily as optimization in
curvature-transition space rather than position space.
\end{remark}

% ============================================================
\section{Canonical Interpretation}
% ============================================================

\begin{proposition}
Transition curves are fundamentally curvature-transition laws.
\end{proposition}

\begin{proposition}
Within AIM, transition geometry is interpreted intrinsically within
curvature space.
\end{proposition}

\begin{proposition}
Sparse transition models therefore represent compact transition
operators rather than explicit geometric curves.
\end{proposition}

% ============================================================
\section{Connection to AIM}
% ============================================================

\begin{summary}
Curvature transitions form the dynamic geometric core of AIM.

Within this framework:
\begin{itemize}
\item transitions are interpreted as curvature evolution processes,
\item sparse models become transition operators,
\item normalized transition space enables reusable families,
\item and runtime systems reconstruct geometry dynamically from
curvature structure.
\end{itemize}

This interpretation unifies transition theory, realization behaviour,
runtime evaluation, and optimization within one intrinsic curvature
framework.
\end{summary}
